\documentclass[11pt]{article}
\input{preamble}
\begin{document}

\begin{center}
{\footnotesize\color{acc}\textsc{ORF 526 \quad$\cdot$\quad Probability for Modern Machine Learning \quad$\cdot$\quad Fall 2026}}\\[0.5em]
{\LARGE\bfseries Lectures 1 and 2: Gaussians}\\[0.35em]
{\color{grey} Definition, linear structure, and conditioning as projection}
\end{center}
\vspace{0.4em}
\hrule height 0.5pt
\vspace{0.6em}

\begin{roadmap}
{\color{acc}\textsc{outline}}\ \ $\varphi_X=\exp(\text{quadratic})$\rmto$AX+b$ is Gaussian\rmto uncorrelated $\Rightarrow$ independent\rmto conditioning is orthogonal projection\rmto ridge regression\\[0.35em]
{\color{acc}\textsc{how to read}}\ \ Sections marked \textcolor{acc}{$\ast$} are not lectured. They are complete, they are examinable, and they carry most of the exercises. Exercises are typed \textsc{a} (did you follow the lecture), \textsc{b} (did you follow the argument) and \textsc{c} (one idea, short once you have it). Each week's quiz has one of each.
\end{roadmap}

\section{Gaussians, and their linear structure}

\subsection{What is a Gaussian?}

Here are three answers, each of which is taught somewhere.

\begin{enumerate}[label=(\alph*),leftmargin=1.6em,itemsep=0.15em]
\item $X$ has density proportional to $\exp(-\tfrac12 (x-\mu)\tr\Sigma^{-1}(x-\mu))$.
\item $X=AZ+\mu$, where $Z$ has independent standard normal coordinates.
\item Every linear functional $\ip{v}{X}$ is a one-dimensional Gaussian.
\end{enumerate}

They are not equivalent. Answer (a) requires $\Sigma$ to be invertible, so it refuses to call $X=(g,g)$ Gaussian --- a vector obtained from a standard normal by a linear map, all of whose linear functionals are Gaussian. A definition that has to apologize for its most obvious example is the wrong definition.

\begin{question}
\textbf{Q.} What should we ask of a definition of the Gaussian?\\[0.25em]
That the class be closed under the operations we care about --- and the operation we care about is applying a linear map.
\end{question}

We take a fourth answer, which implies the others and costs nothing at the degenerate boundary.

\subsection{What we import}

\begin{definition}[Characteristic function]
For a random vector $X$ in $\R^n$,
\[
\varphi_X(t)\;=\;\E\big[e^{i\ip{t}{X}}\big],\qquad t\in\R^n .
\]
\end{definition}

It is defined for every $X$ with no integrability assumption, because $|e^{i\theta}|=1$. Immediately: $\varphi_X(0)=1$, $|\varphi_X(t)|\le 1$, and $\varphi_X(-t)=\overline{\varphi_X(t)}$.

\begin{imported}
\textbf{Imported: Fourier uniqueness and inversion.} If $\varphi_X=\varphi_Y$ then $X\dto Y$. If moreover $\varphi_X\in L^1(\R^n)$, then $X$ has a bounded continuous density
\[
p(x)\;=\;\frac{1}{(2\pi)^n}\int_{\R^n} e^{-i\ip{t}{x}}\varphi_X(t)\,dt .
\]
This is standard Fourier analysis and the only thing we do not prove.
\end{imported}

\subsection{The definition}

\begin{definition}[Gaussian vector]\label{def:gauss}
Let $\mu\in\R^n$ and let $\Sigma$ be symmetric $n\times n$. We say $X\sim\Nor(\mu,\Sigma)$ if
\[
\boxed{\ \varphi_X(t)\;=\;\exp\Big(i\ip{t}{\mu}-\tfrac12\,t\tr\Sigma t\Big)\ }
\]
for all $t\in\R^n$.
\end{definition}

No inverse of $\Sigma$ appears, so nothing prevents $\Sigma$ from being singular. The definition describes a law rather than constructing one, so we owe an existence proof; we give it in \S1.5, where it will take a single line.

\subsection{Linear structure}

\begin{theorem}[Linear images]\label{thm:lin}
If $X\sim\Nor(\mu,\Sigma)$ and $A\in\R^{m\times n}$, $b\in\R^m$, then
\[
AX+b\;\sim\;\Nor\big(A\mu+b,\;A\Sigma A\tr\big).
\]
\end{theorem}

\begin{proof}
$\varphi_{AX+b}(t)=e^{i\ip{t}{b}}\E[e^{i\ip{A\tr t}{X}}]=e^{i\ip{t}{b}}\varphi_X(A\tr t)$. Substituting Definition~\ref{def:gauss} and using $\ip{A\tr t}{\mu}=\ip{t}{A\mu}$ gives $\exp\big(i\ip{t}{A\mu+b}-\tfrac12 t\tr A\Sigma A\tr t\big)$.
\end{proof}

\begin{remark}
It is $A\Sigma A\tr$, not $A\tr\Sigma A$.
\end{remark}

\begin{corollary}[Answer (c), recovered]\label{cor:1d}
$X\sim\Nor(\mu,\Sigma)$ if and only if $\ip{v}{X}\sim\Nor(\ip{v}{\mu},\,v\tr\Sigma v)$ for every $v\in\R^n$.
\end{corollary}

\begin{proof}
Forward: take $A=v\tr$ in Theorem~\ref{thm:lin}. Conversely, evaluating the one-dimensional transform of $\ip{v}{X}$ at $1$ gives $\varphi_X(v)=\exp(i\ip{v}{\mu}-\tfrac12 v\tr\Sigma v)$, which is Definition~\ref{def:gauss}.
\end{proof}

\subsection{The standard normal, and existence}

One computation is still missing: we have never exhibited a single Gaussian.

\begin{lemma}\label{lem:phig}
If $g$ is standard normal on $\R$ then $\E[e^{i\lambda g}]=e^{-\lambda^2/2}$; equivalently $g\sim\Nor(0,1)$.
\end{lemma}

\begin{proof}
Set $f(\lambda)=\E[e^{i\lambda g}]$. Differentiating under the integral,
\[
f'(\lambda)\;=\;\frac{1}{\sqrt{2\pi}}\int ix\,e^{i\lambda x}e^{-x^2/2}\,dx .
\]
The key observation is that $x\,e^{-x^2/2}=-\frac{d}{dx}e^{-x^2/2}$, so integrating by parts moves the derivative onto the exponential:
\[
f'(\lambda)\;=\;\frac{i}{\sqrt{2\pi}}\int \Big(\frac{d}{dx}e^{i\lambda x}\Big)e^{-x^2/2}\,dx\;=\;i\cdot i\lambda\, f(\lambda)\;=\;-\lambda f(\lambda).
\]
With $f(0)=1$, $f(\lambda)=e^{-\lambda^2/2}$.
\end{proof}

\begin{remark}
The proof used $\E[g\,h(g)]=\E[h'(g)]$ with $h(x)=e^{i\lambda x}$: multiplication by $g$ acts like differentiation. This identity returns in Lecture 4 under its own name.
\end{remark}

If $Z=(g_1,\dots,g_n)$ has independent standard normal coordinates then $\varphi_Z(t)=\prod_j e^{-t_j^2/2}=e^{-|t|^2/2}$, so $Z\sim\Nor(0,I_n)$. Existence is now immediate.

\begin{proposition}[Well-posedness]\label{prop:wp}
A random vector satisfying Definition~\ref{def:gauss} exists if and only if $\Sigma\succeq0$, and is then unique in law. If $\Sigma=AA\tr$, one may take $X=AZ+\mu$.
\end{proposition}

\begin{proof}
If $\Sigma\succeq0$, factor $\Sigma=AA\tr$. Since $Z\sim\Nor(0,I)$, Theorem~\ref{thm:lin} gives $AZ+\mu\sim\Nor(\mu,AA\tr)=\Nor(\mu,\Sigma)$. Uniqueness in law is Fourier uniqueness.

Conversely, suppose $v\tr\Sigma v=-c$ with $c>0$. Along $t=sv$ the claimed transform has modulus $e^{s^2c/2}\to\infty$, contradicting $|\varphi_X|\le1$.
\end{proof}

The family of Gaussian laws on $\R^n$ is therefore indexed exactly by $\R^n\times\{\Sigma\succeq0\}$. Answer (b) has become a construction inside a proof rather than a rival definition.

\begin{corollary}[Rotational invariance]\label{cor:rot}
If $Z\sim\Nor(0,I_n)$ and $O$ is orthogonal then $OZ\dto Z$.
\end{corollary}

\begin{proof}
$\varphi_{OZ}(t)=\varphi_Z(O\tr t)=e^{-|O\tr t|^2/2}=e^{-|t|^2/2}$.
\end{proof}

Invariance is what singles the Gaussian out among product measures. In November, invariance under orthogonal conjugation will single out a particular matrix ensemble in the same way.

\begin{example}[The degenerate case]
Take $n=2$, $\mu=0$, $\Sigma=\begin{psmallmatrix}1&1\\1&1\end{psmallmatrix}$, which is positive semidefinite of rank one. By Proposition~\ref{prop:wp} the law $\Nor(0,\Sigma)$ exists, and $A=\begin{psmallmatrix}1\\1\end{psmallmatrix}$ realizes it as $X=(g,g)$. It has no density --- it is supported on a line --- and Definition~\ref{def:gauss} never notices.
\end{example}

\subsection{Uncorrelated versus independent}

For general random variables zero covariance says very little. For jointly Gaussian ones it says everything.

\begin{theorem}\label{thm:indep}
Let $(X,Y)$ be jointly Gaussian in $\R^n\times\R^m$ --- the concatenated vector is Gaussian --- with $\Cov(X,Y)=0$. Then $X\indep Y$.
\end{theorem}

\begin{proof}
Write $t=(s,u)$. The covariance of the concatenated vector is block diagonal, so $t\tr\Sigma t=s\tr\Sigma_{XX}s+u\tr\Sigma_{YY}u$ and $\ip{t}{\mu}=\ip{s}{\mu_X}+\ip{u}{\mu_Y}$. Hence
\[
\varphi_{(X,Y)}(s,u)\;=\;\varphi_X(s)\,\varphi_Y(u),
\]
which is the characteristic function of the product law $\mathcal{L}(X)\otimes\mathcal{L}(Y)$. By Fourier uniqueness the joint law is that product.
\end{proof}

The hypothesis that fails in practice is joint Gaussianity, not Gaussianity of the marginals.

\begin{example}[Uncorrelated, standard normal, not independent]\label{ex:fail}
Let $g\sim\Nor(0,1)$ and let $\varepsilon$ be an independent random sign; set $X=g$, $Y=\varepsilon g$. Conditionally on $\varepsilon=\pm1$ we have $Y=\pm g\dto g$, so both marginals are standard normal. They are uncorrelated, since $\E[XY]=\E[\varepsilon]\E[g^2]=0$. And $|X|=|Y|$ always, so $X$ determines $Y$ up to a sign.

Conditioning on $\varepsilon$ and applying Lemma~\ref{lem:phig},
\[
\varphi_{(X,Y)}(s,u)\;=\;\tfrac12 e^{-(s+u)^2/2}+\tfrac12 e^{-(s-u)^2/2},
\]
a sum of two Gaussian transforms and not $\exp(\text{quadratic})$. The pair fails Definition~\ref{def:gauss} even though each coordinate satisfies it.
\end{example}

\subsection{What we have not proved}

\textbf{Maximum entropy.} Among all laws on $\R^n$ with mean $\mu$ and covariance $\Sigma\succ0$, the Gaussian uniquely maximizes differential entropy: it is the least committal law consistent with a given second moment. Section~\ref{sec:entropy} proves it, and Lecture 3 uses it.

\textbf{The moment generating function.} The computation of Lemma~\ref{lem:phig} with a real exponent gives $\E[e^{\lambda g}]=e^{\lambda^2/2}$, from which Gaussian tail bounds follow by Chernoff's method. In Module II this becomes the engine of large deviations.

\starsec{Entropy, and why the Gaussian maximizes it}\label{sec:entropy}

\emph{Not lectured. Exercises due at the start of Lecture 2.}

Before entropy, surprise. If an event of probability $p$ occurs, the surprise $s(p)$ ought to satisfy $s(1)=0$ and $s(pq)=s(p)+s(q)$. Substituting $p=e^{-u}$ turns the second into $g(u+v)=g(u)+g(v)$ for $g(u)=s(e^{-u})$, whose only continuous solutions are linear. So $s(p)=-c\log p$, the base of the logarithm is a choice of units, and we take $c=1$.

The \textbf{entropy} of a law on a finite set is the surprise expected before looking,
\[
H(p)=\E\big[-\log p(X)\big]=-\sum_xp(x)\log p(x).
\]

\begin{exercise}[\ty{A}]
Compute $H$ for a coin of bias $q$; sketch it on $q\in[0,1]$ and find where it is largest. Compute $H$ for the uniform law on $n$ points.
\end{exercise}

\begin{exercise}[\ty{B}]
Show $H(p)\le\log n$ for every law on $n$ points, with equality only for the uniform one. \emph{Apply Jensen to $\E[\log(1/p(X))]$ and notice what $\E[1/p(X)]$ is.}
\end{exercise}

That is the pattern we are after: \emph{fix a constraint, and entropy is maximized by the least committal law satisfying it.} With no constraint on a finite set that is the uniform law; the Gaussian will answer the same question with a different constraint.

For a density $p$ on $\R^n$ the same formula defines the \textbf{differential entropy} $h(p)=-\int p\log p$. It is not the discrete entropy in disguise.

\begin{exercise}[\ty{A}]
Compute $h$ for the uniform law on $[0,a]$, the exponential of rate $\lambda$, and $\Nor(0,\sigma^2)$ --- you should find $\tfrac12\log(2\pi e\sigma^2)$ for the last. Exhibit a law with \emph{negative} differential entropy, and say why no discrete entropy can be negative.
\end{exercise}

\begin{exercise}[\ty{A}]
Show $h(aX)=h(X)+\log|a|$ for $a\ne0$, so differential entropy is not invariant under a change of units: ``the entropy of a length'' is not a well-defined number.
\end{exercise}

The cure is to measure entropy \emph{relative} to a reference. For densities $p,q$, the \textbf{relative entropy} is $\KL{p}{q}=\int p\log(p/q)$ --- the extra surprise from expecting $q$ when the truth is $p$, so it should never be negative.

\begin{exercise}[\ty{B}]
Prove \textbf{Gibbs' inequality}: $\KL{p}{q}\ge0$, with equality iff $p=q$. \emph{Jensen applied to $-\log$, exactly as before.} Check that $\KL{p}{q}$ is unchanged by rescaling $X$ --- unlike $h$.
\end{exercise}

Now fix $\mu$ and $\Sigma\succ0$, let $q$ be the Gaussian density, and let $p$ be any density with that mean and covariance. Expanding $0\le\KL{p}{q}$ gives $0\le-h(p)-\int p\log q$, and everything hinges on the second term.

\begin{exercise}[\ty{B}]
Write out $\log q$ and observe it is a \emph{quadratic polynomial} in $x$. Deduce that $\int p\log q$ depends on $p$ only through its mean and covariance, hence equals $\int q\log q=-h(q)$. Conclude $h(p)\le h(q)$ with equality only when $p=q$, and check $h(q)=\tfrac12\log\big((2\pi e)^n\det\Sigma\big)$.
\end{exercise}

Note how little the Gaussian had to do: everything turned on $\log q$ being quadratic, so that integrating it against $p$ could see only $p$'s first two moments. \emph{A quadratic is exactly what a covariance constraint measures.}

\begin{exercise}[\ty{C}]
Let $X,Y$ be i.i.d., \emph{symmetric}, each with a density. Show
\[
h\Big(\frac{X+Y}{\sqrt2}\Big)\;\ge\;h(X).
\]
\emph{You may use subadditivity $h(U,V)\le h(U)+h(V)$, which is the previous exercise applied to a joint density against the product of its marginals.} Lecture 3 opens with this inequality.
\end{exercise}

\begin{exercise}[\ty{C}]
Fix $S\subseteq\R^n$, functions $T_1,\dots,T_k$ on $S$ and numbers $c_1,\dots,c_k$. \textbf{Assume} that for some $\lambda_1,\dots,\lambda_k$ the function $\exp(\sum_m\lambda_mT_m)$ has finite integral $Z$ over $S$, so that $q=Z^{-1}\exp(\sum_m\lambda_mT_m)$ is a density, and that $\int qT_m=c_m$ for each $m$. Show $q$ uniquely maximizes $h$ among densities on $S$ with those moments.

Then reread the paragraph above --- \emph{which property of the Gaussian did that argument use?} Recover as instances the uniform law on $[0,a]$, the exponential at fixed mean on $[0,\infty)$, and $\Nor(\mu,\Sigma)$ at fixed mean and covariance. Finally, show the assumption does real work: on $S=\R$ with the mean alone constrained no $\lambda$ makes $e^{\lambda x}$ integrable, and there is no maximizer --- exhibit densities of mean $0$ with $h$ arbitrarily large.
\end{exercise}

\section{Conditioning as projection}

\begin{question}
\textbf{Q.} We observe $Y$. What is our best guess for $X$?\\[0.35em]
Before answering we must say what \emph{best} means.
\end{question}

Fix jointly Gaussian $(X,Y)$ in $\R\times\R^m$, both centered for readability.

\subsection{What would count as an answer}

Measure a guess by mean squared error. A predictor is any function of the observation, and we want the one minimizing $\E[(X-f(Y))^2]$.

\begin{definition}[Conditional expectation]\label{def:ce}
$\E[X\mid Y]$ is the minimizer of $\E[(X-f(Y))^2]$ over all measurable $f$ with $\E f(Y)^2<\infty$.
\end{definition}

One could instead define $\E[X\mid Y]$ as the orthogonal projection of $X$ onto the span of the coordinates of $Y$. That would make Theorem~\ref{thm:proj} a tautology. Under Definition~\ref{def:ce} the competitors are all functions of $Y$, including nonlinear ones.

\subsection{The right inner product}

On the linear span of the coordinates of $X$ and $Y$ --- a finite-dimensional space of random variables --- put
\[
\ip{U}{V}\;:=\;\E[UV].
\]
Inner product is covariance; orthogonal is uncorrelated. Let $H_Y=\operatorname{span}\{Y_1,\dots,Y_m\}$ and let $\Pi$ be orthogonal projection onto $H_Y$.

\begin{lemma}[Best linear predictor]\label{lem:blp}
$\Pi X=\Sigma_{XY}\Sigma_{YY}^{-1}Y$, with $\Sigma_{XY}=\E[XY\tr]$ and $\Sigma_{YY}=\E[YY\tr]$. If $\Sigma_{YY}$ is singular the same formula holds with a pseudo-inverse, and $\Pi X$ remains unique.
\end{lemma}

\begin{proof}
Write $\Pi X=\sum_j c_jY_j$. Orthogonality of the residual to each $Y_k$ reads $\E[(X-\sum_j c_jY_j)Y_k]=0$, that is $\Sigma_{YY}c=\Sigma_{YX}$.
\end{proof}

Lemma~\ref{lem:blp} holds for any square-integrable pair.

\subsection{The theorem}

\begin{theorem}[Conditioning is projection]\label{thm:proj}
Let $(X,Y)$ be jointly Gaussian and centered. Then
\[
\E[X\mid Y]\;=\;\Pi X\;=\;\Sigma_{XY}\Sigma_{YY}^{-1}Y .
\]
\end{theorem}

\begin{proof}
Let $R=X-\Pi X$. By construction $\Cov(R,Y)=\Sigma_{XY}-\Sigma_{XY}\Sigma_{YY}^{-1}\Sigma_{YY}=0$, and $(R,Y)$ is a linear image of $(X,Y)$, hence jointly Gaussian by Theorem~\ref{thm:lin}. So Theorem~\ref{thm:indep} gives $R\indep Y$.

For any competitor $f$,
\[
\E\big[(X-f(Y))^2\big]=\E[R^2]+\E\big[(\Pi X-f(Y))^2\big]+2\,\E\big[R\,(\Pi X-f(Y))\big].
\]
The cross term vanishes: $\Pi X-f(Y)$ is a function of $Y$, and $R$ is independent of $Y$ with mean zero, so the expectation factors. Hence $\E[(X-f(Y))^2]\ge\E[R^2]$, with equality exactly when $f(Y)=\Pi X$ almost surely.
\end{proof}

\begin{remark}
Uncorrelatedness of $R$ and $Y$ does not suffice to kill the cross term; independence does, and independence comes from Theorem~\ref{thm:indep}.
\end{remark}

\subsection{The conditional law}

Since $X=\Pi X+R$ with $R\indep Y$, conditioning on $Y=y$ leaves $R$ untouched.

\begin{corollary}[Schur complement]\label{cor:schur}
For jointly Gaussian $(X,Y)$,
\[
X\mid Y=y\;\sim\;\Nor\Big(\mu_X+\Sigma_{XY}\Sigma_{YY}^{-1}(y-\mu_Y),\ \ \Sigma_{XX}-\Sigma_{XY}\Sigma_{YY}^{-1}\Sigma_{YX}\Big).
\]
\end{corollary}

\begin{remark}
The conditional covariance does not depend on $y$.
\end{remark}

\begin{remark}
Conditioning on $\{Y=y\}$ conditions on an event of probability zero. The decomposition $X=\Pi X+R$ with $R\indep Y$ constructs a conditional law directly, so Corollary~\ref{cor:schur} is a definition justified by Theorem~\ref{thm:proj}.
\end{remark}

\subsection{Bayesian linear models and ridge regression}

Take
\[
y \;=\; Hw+\sigma\xi,\qquad w\sim\Nor(0,\tau^2 I_p),\quad \xi\sim\Nor(0,I_n)\ \text{independent of } w .
\]
The pair $(w,y)$ is a linear image of $(w,\xi)$, hence jointly Gaussian by Theorem~\ref{thm:lin}. Reading off blocks,
\[
\Sigma_{ww}=\tau^2 I,\qquad \Sigma_{wy}=\tau^2H\tr,\qquad \Sigma_{yy}=\tau^2HH\tr+\sigma^2I,
\]
and Corollary~\ref{cor:schur} gives the posterior. Simplifying with $\tau^2H\tr(\tau^2HH\tr+\sigma^2I)^{-1}=(H\tr H+\lambda I)^{-1}H\tr$:

\begin{proposition}\label{prop:ridge}
With $\lambda=\sigma^2/\tau^2$,
\[
\E[w\mid y]=\big(H\tr H+\lambda I\big)^{-1}H\tr y,
\qquad
\Cov(w\mid y)=\sigma^2\big(H\tr H+\lambda I\big)^{-1}.
\]
\end{proposition}

The posterior mean is ridge regression, with regularization parameter the ratio of noise variance to prior variance. Three claims are now visible at once: \emph{a loss function is a noise model, a regularizer is a prior, and fitting is conditioning.}

\emph{The exercises below are due at the start of Lecture 3.}

\begin{exercise}[\ty{A}]
Show $-\log p(y\mid w)=\frac{1}{2\sigma^2}\|y-Hw\|^2+c$ with $c$ free of $w$, so least squares is maximum likelihood. Say in one sentence what someone assumes about the world when they choose squared error, and whether they usually know it.
\end{exercise}

\begin{exercise}[\ty{A}]
Show the maximizer of the posterior density is $\widehat w_\lambda=\arg\min_w\|y-Hw\|^2+\lambda\|w\|^2$. What does a large $\lambda$ say about the data relative to the prior, and what are the limits $\tau\to0$ and $\tau\to\infty$?
\end{exercise}

\begin{exercise}[\ty{B}]
Fill in the derivation above: compute the three blocks of the joint covariance and apply Corollary~\ref{cor:schur}. Compare its length with completing the square in the density. Then note that $\Cov(w\mid y)$ does not depend on $y$, and say what that means operationally --- the data came out surprising, so by how much should you revise your uncertainty?
\end{exercise}

The posterior mean is linear in $y$, so by Theorem~\ref{thm:proj} the best predictor of $w$ happens to be linear. That is special.

\begin{exercise}[\ty{B}]
Let $Y\sim\Nor(0,1)$ and $X=Y^2-1$. Show $\Cov(X,Y)=0$, so the best \emph{linear} predictor of $X$ is the constant $0$ with mean squared error $2$. Exhibit a predictor with error $0$. Which hypothesis of Theorem~\ref{thm:proj} fails?
\end{exercise}

\starsec{The two forms of the ridge estimator}

\emph{Not lectured.} Proposition~\ref{prop:ridge} inverts a $p\times p$ matrix; the derivation that produced it inverted an $n\times n$ one. Both are correct.

\begin{exercise}[\ty{C}]
Show that for $\lambda>0$
\[
\big(H\tr H+\lambda I_p\big)^{-1}H\tr\;=\;H\tr\big(HH\tr+\lambda I_n\big)^{-1},
\]
and deduce that ridge regression with $p\gg n$ costs $O(n^3)$ rather than $O(p^3)$. Then take $\tau\to\infty$ with $n<p$ and $H$ of full row rank: show the posterior mean converges to $H\tr(HH\tr)^{-1}y$, the \textbf{minimum-norm} solution of $Hw=y$.

\emph{Ordinary least squares does not exist in this regime and the posterior does. The estimator you have just found is the one that interpolates the data, and Module II explains why its test error behaves as it does.}
\end{exercise}

\end{document}
